\scanpage{28}
\mspage{2}
les surfaces suivantes :

a) le plan projectif réel ($\simeq$ Sphère conforme/antipod.) et le plan projectif réel privé d'un point (dont les rev. universels hol. resp. $\simeq\mathbb P^1_{\mathbb C}$ et $\simeq G_{m\,\mathbb C}$).

b) Les surfaces conformes associées aux courbes \deleted{elliptiques}\inserted{lisses géom. connexes} réelles $E$\inserted{de genre $1$} telles que $E(\mathbb R)=\varnothing$, en posant \deleted{$\unsure{E_{\mathbb C}}$}
\[
X=E(\mathbb C)/\text{conj. complexe}.
\]
On trouve que les $E$ en question, à isom. près, sont les $E_a^!$, $a\in\mathbb R^{*+}$\ldots

\rule{\linewidth}{0.3pt}

\uline{Question des \inserted{\unsure{la}} métriques riemanniennes}\\
\uline{canoniques sur une surface conforme}

(sans bord pour commencer)

\textcircled{1} Cas hyperbolique, Il y a une métrique canonique, à courbure constante de valeur $-1$, déduite par descente du revêtement universel $\widetilde X$.

\uline{Métrique hyperbolique}
% archive foliation: 30
